\section{Introduction}
\label{sec:intro}


Deep learning is famously a black-box learning method, the most powerful, most inscrutable, and now most technologically important member of the machine learning pantheon.
Properly trained, neural networks learn to perform a wide array of tasks with superhuman performance, but we have no unified scientific framework that explains why or how.
Motivated by both scientific curiosity and the promise of practical engineering benefit, the effort to put rigorous mathematical and scientific backing behind this applied discipline has spanned decades.
Despite some progress, however, our understanding remains primitive: neural networks are still trained using methods discovered largely through trial and error rather than first principles, and theory plays little role in the day-to-day practice of deep learning.
The challenge has only compounded as practice has advanced, and in the era of large language models and diffusion models, the mysteries are arguably deeper than they were one or two decades ago.
Will we ever understand?

\vspace{2mm}
\begin{tcolorbox}
[sharp corners, colback=white, colframe=black]
    \textbf{
    \hspace{-1.5mm}
    This paper makes the case that, yes, there will be a scientific theory of deep learning; that we can see pieces of this theory starting to emerge; and that this theory will take the form of a \textit{mechanics} of the learning process.
    }
\end{tcolorbox}
\vspace{2mm}


The questions driving deep learning theory have changed over time, and to understand where the field is going, it is useful to first look back at how we got here.
Deep learning theory is as old as machine learning itself, with roots in the McCulloch--Pitts neuron and the perceptron in the middle of the last century.
The earliest theoretical questions in machine learning were about expressivity: what functions can simple models represent, and how can they be learned from data?
As learning came to be understood as a statistical problem, and simple learning systems found practical success, the theoretical focus shifted to ask: when does learning from finite samples generalize?
This gave rise to \textit{classical learning theory}, including statistical and computational/PAC learning theory.
Paired with \textit{classical optimization theory,} these frameworks gave clean end-to-end guarantees of the optimization and generalization of simple learning systems.
In parallel, a classical tradition of the \textit{statistical physics of machine learning} developed satisfying theories of the average-case behavior of simple models.



While these classical theories built a strong foundation for understanding learning, the rise of deep learning through multilayer networks, backpropagation, and increasing scale in both data and compute exposed limitations in their explanatory power.
Neural networks are complex, nonconvex, and overparameterized (in contrast with the simple, convex, parsimonious models for which classical learning theory excels), and they optimize and generalize better than these classical approaches can guarantee or explain.
Furthermore, it became clear that neural networks were not merely fitting data or achieving low training error, they were learning structured internal representations and displaying striking regularities across tasks and scales.
The classical questions of performance and efficiency remained important, but answering them would first require understanding a new host of phenomena shaped both by the dynamics of neural networks through training and the structure of the data they are trained on.


This marked a transition in which deep learning theory changed in character from a largely \textit{mathematical} study of what is possible to a truly \textit{scientific} effort to describe, explain, and ultimately predict the behavior of complex empirical systems.
New scientific endeavors often start with an empirical tension in which nature presents something interesting we cannot predict or explain with existing tools, and although neural networks are artificial computational systems, this same scientific tension is present here.
We should thus approach this task as scientists, embracing empirics, seeking unifying principles, and identifying recurring motifs.
We should also expect the path forward to look more like the development of a scientific field than the development of a mathematical one.

The purpose of this paper is to convince the reader that this scientific tension is gradually giving way to a scientific theory which resolves it.
In \cref{sec:reasons_why}, we pull together major strands of ongoing research and identify five lines of evidence that such a theory is emerging:

\begin{enumerate}
    \item there are a growing number of \textit{analytically solvable settings} in which learning is fully captured by simple mathematics, including but not limited to deep linear networks and kernel methods (\cref{sec:reason-dynamics});
    \item there exist \textit{useful limits}, including limits of infinite width and depth, that provide insight into fundamental learning behaviors (\cref{sec:reason-limits});
    \item in many cases, \textit{simple empirical laws} suffice to capture meaningful macroscopic statistics, including test-time performance and loss-landscape sharpness (\cref{sec:reason-measurable});
    \item many of the hyperparameters governing optimization can be \textit{disentangled and understood}, leaving behind a simpler effective dynamical system (\cref{sec:reason-hps}); and
    \item as applied deep learning has scaled up and converged on best practices, \textit{universal phenomena} have increasingly appeared across settings and tasks (\cref{sec:reason-universal-behavior}).
\end{enumerate}

These lines of research broadly share several overarching characteristics: they are concerned with the \textit{dynamics of the training process}; they primarily seek to describe \textit{coarse aggregate statistics} of learning; and they emphasize \textit{accurate average-case predictions} over rigorous worst-case bounds.
In this sense, the emerging scientific theory of neural networks appears to have much in common with theories in physics such as classical mechanics, continuum mechanics, statistical mechanics, and quantum mechanics.
We argue that this emerging theory is best understood as a \textit{mechanics of learning}.



\subsection{What's in a mechanics?}

Mechanics is the branch of physics studying how forces acting on objects determine their movement through space and time.
Neural network learning can be thought of in this way: much as an object moves continuously through physical space, learning involves a model moving through parameter space via discrete updates.
In the physical sciences, forces come from interactions between components of a system.
Similarly, the process of deep learning is shaped by interactions between the parameters, dataset, task, and learning rule.
In physics, these forces are mediated by fields; in deep learning, they are mediated by gradients.
In physics, systems settle into equilibria at local minima of a potential determined by internal interactions and external constraints; analogously, neural networks converge to local minima of a loss landscape shaped by their architecture and training data.
While the systems under study are very different, since the key problems of both are essentially about movement and interaction, we might expect some features of the resulting sciences to be shared.

These analogies are not just speculation: we can see these similarities reflected in the lines of research listed above.
All branches of mechanics (and especially classical mechanics) develop a library of analytically solvable settings to gain intuition; so too does learning mechanics.
All branches of mechanics use limits as simplifying tools; so too does learning mechanics.
Continuum and statistical mechanics, the branches which most directly deal with large numbers of interacting components, describe zoomed-out summary statistics rather than the motion of every particle; this has also proven a useful approach in dealing with the complexity of deep learning.
Every physical system has one or more system parameters (characteristic scales, coupling constants, etc.) affecting its behavior, and some techniques for treating these are essentially the same as those used to study hyperparameters in deep learning.
Finally, physics is full of cases in which the same phenomena show up in very different settings, and similarly we see universal behavior emerging across deep learning systems.

All considered, the emerging science shares deep similarities with established branches of mechanics.
By analogy to classical, continuum, statistical, and quantum mechanics, we suggest the intended theory be called \textit{learning mechanics}.




















\paragraph{Seven desiderata for learning mechanics.}


We should be clear at the outset what we want from a mechanics of learning.
Assessing how mature branches of mechanics were motivated, developed, and succeeded, we can see what sort of goals to aim for.
Here are seven desiderata for this research program:

\begin{enumerate}[itemsep=1ex]
    \item
    Learning mechanics should be \textit{fundamental}, proceeding logically from a first-principles description of neural network training.
    Interim assumptions about network weights, dynamics, and performance will be useful tools, but they should ultimately be explained from first principles.
    \item
    Learning mechanics should be \textit{mathematical}, making unambiguous quantitative statements about important properties of neural networks.
    No mechanics is a qualitative science; neither will be learning mechanics.
    \item
    Learning mechanics should be \textit{predictive}, making claims supported by simple, repeatable empirical measurements.
    We have excellent experimental control of our system, and every major development should be unambiguously verified in experiment.
    \item Learning mechanics should be \textit{comprehensive}, describing aspects of neural networks' training process, hidden representations, and final weights in a single picture.
    It is worth emphasizing that this theory will not --- and should not --- aim to describe everything.
    A map at the full resolution of the world would be the size of the world and thus of little use.
    What we seek instead is a theory that operates at the right level of resolution --- one that sacrifices detail in favor of insight.
    \item Learning mechanics should be \textit{intuitive}, being simple, illuminating, and satisfying in its demystification of deep learning.
    Like physics, learning mechanics should strive for simple insight over technical complexity.
    \item Learning mechanics should be \emph{useful}, serving as the scientific foundation for applied deep learning as physics does for other forms of engineering.
    Concrete goals should include greatly reducing the need for hyperparameter tuning, giving predictive tools for dataset design, and providing rigorous foundations for AI safety work.
    \item Finally, learning mechanics should be \textit{humble}, being solid in what it describes and explicit about what it cannot.
    Every branch of physical science has a regime of applicability outside of which it breaks down, and these boundaries are taught together with the science so that it may be used reliably.
    We anticipate the mechanics of learning applicable to realistic deep learning will break down in many small-scale, handcrafted, or otherwise special cases, and this is the price we will pay for the right simple picture in the regimes we care about.
\end{enumerate}

A mechanics of learning with these virtues --- one that is fundamental, mathematical, predictive, comprehensive, intuitive, useful, and humble --- would be transformative, paradigm-setting.
We expect such a theory would resolve important open questions that have long remained out of reach, as we discuss in \cref{sec:open_dirs}.

\subsection{Why learning mechanics matters}

Building learning mechanics will not be easy.
It will require sustained effort, both intellectual and institutional.
It is therefore worth being clear about why such a project matters.
The reasons to seek a mechanics of learning fall into three broad categories: \textit{scientific}, \textit{practical}, and \textit{safety-related}.

The scientific reasons concern what such a theory could teach us about intelligence and the natural world.
The striking engineering success of large neural networks suggests that they exploit deep principles of learning and representation that we do not yet understand.
This has historical precedent: technology has often preceded scientific theory, as was the case with steam engines' role in motivating thermodynamics, which went on to explain much more than engine efficiency.
A similar story played out in flight: the development of airplanes through trial and error and inspiration from the natural world helped motivate aerodynamic theory, which in turn enabled both better aircraft design and a deeper understanding of how birds themselves fly.
In our case, the principles that govern learning in artificial neural networks may also shed light on our own biological intelligence, with potentially important implications for neuroscience and cognitive science.

The practical reasons concern the design and development of real-world AI systems.
A mature theory of deep learning could guide model design, optimization, scaling, and deployment, replacing trial and error with more reliable principles.
Theory has already begun to play this role in a limited but growing number of cases, including empirical scaling laws (\cref{sec:reason-measurable}), mathematical prescriptions for hyperparameter scaling (\cref{sec:reason-hps}), and theoretically-motivated optimizers and methods for data attribution (\cref{sec:reasons_for_skepticism}).
A deeper, more complete theory will give more such guidance and make it sharper and more predictive.

The safety reasons concern our ability to describe, characterize, and govern increasingly powerful AI systems.
Some form of regulation will likely be necessary, but it is difficult to regulate a technology that we cannot clearly describe.
A theory that identifies the relevant variables, mechanisms, and organizing principles of large models could help provide the clarity needed for reliability, oversight, and control.
One avenue by which fundamental theory might aid in AI safety is by supporting mechanistic interpretability, a point to which we return in \cref{sec:learning_mech_and_mechinterp}.


\subsection{Plan for this paper}

This paper is structured as follows.
In \cref{sec:reasons_why}, we present our five lines of evidence that a scientific theory of deep learning is beginning to emerge.
We motivate each line of evidence with an intuitive explanation and highlight examples of research successes that illustrate the underlying principle.
In \cref{sec:learning_mech_and_mechinterp}, we discuss the relationship between learning mechanics and other perspectives on the science of deep learning, including a possible symbiotic relationship between learning mechanics and mechanistic interpretability.
In \cref{sec:reasons_for_skepticism}, we review and address common arguments that fundamental theory will not be possible.
In \cref{sec:open_dirs}, we give a portrait of ten important open directions in learning mechanics, from predicting scaling laws to eliminating hyperparameters, where we expect to see major progress in the coming years.
Finally, in \cref{sec:how_to_get_involved}, we offer some advice for young researchers looking to get involved in this scientific project and extend a hand with some introductory resources.

We write this paper for a broad audience.
We hope the \textit{veteran scientist of deep learning} will find something valuable in our synthesis of useful approaches and results, and feel galvanized by our depiction of an emerging science.
We hope to convince the \textit{deep learning practitioner} that theory is on a path to fulfilling its longstanding promise of practical utility and to encourage them to experiment with their systems with an eye for science.
We hope to convince the \textit{AI safety or mechanistic interpretability researcher} that white-box theory is difficult yet possible --- that a first-principles study of dynamics can help put solid foundations beneath their important work, and that our communities should work together (see \cref{sec:learning_mech_and_mechinterp} for our vision of symbiosis).
Lastly, we hope to make it easier for \textit{young students and newcomers to the field} to get involved.
This is an exciting and important area of work, and while it requires some mathematical maturity to get started, it is our belief that the barrier to entry could be much lower.
Various deep intuitions about this science have been percolating inside the theory community for a while, and this paper is an attempt to state them clearly.
We hope to make it easier for folks with the requisite background to quickly get up to speed and contribute.
